\documentclass[11pt,a4paper]{article}

% ===================== Codificación e idioma =====================
\usepackage[utf8]{inputenc}
\usepackage[T1]{fontenc}
\usepackage[spanish,es-tabla,es-noindentfirst]{babel}
\usepackage{lmodern}

% ===================== Color (cargar al inicio) ==================
\usepackage[table]{xcolor}

% Paleta de color corporativa
\definecolor{primario}{HTML}{1F4E79}
\definecolor{secundario}{HTML}{2E75B6}
\definecolor{acento}{HTML}{C00000}
\definecolor{fondoClaro}{HTML}{EAF1F8}
\definecolor{fondoGris}{HTML}{F2F2F2}
\definecolor{verde}{HTML}{2E7D32}
\definecolor{textoGris}{HTML}{333333}

% ===================== Diseño de página ==========================
\usepackage[a4paper,margin=2.3cm,top=2.5cm,bottom=2.5cm]{geometry}
\usepackage{setspace}
\setstretch{1.15}
\setlength{\parskip}{4pt}

% ===================== Tipografía y símbolos =====================
\usepackage{microtype}
\usepackage{textcomp}
\usepackage{fontawesome5}

% ===================== Tablas ====================================
\usepackage{array}
\usepackage{booktabs}
\usepackage{longtable}
\usepackage{tabularx}
\newcolumntype{Y}{>{\raggedright\arraybackslash}X}

% ===================== Encabezados y cajas =======================
\usepackage{fancyhdr}
\usepackage[most]{tcolorbox}

% Caja de nota
\newtcolorbox{cajanota}[1][]{%
  enhanced, breakable, sharp corners,
  colback=fondoClaro, colframe=secundario,
  boxrule=0.4pt, leftrule=3pt,
  fonttitle=\bfseries\color{white},
  coltitle=white, colbacktitle=secundario,
  title={Nota}, #1
}

% Caja de objetivo
\newtcolorbox{cajaobjetivo}[1][]{%
  enhanced, breakable, sharp corners,
  colback=verde!7, colframe=verde,
  boxrule=0.4pt, leftrule=3pt,
  fonttitle=\bfseries\color{white},
  coltitle=white, colbacktitle=verde,
  title={Objetivo}, #1
}

% Caja de alerta
\newtcolorbox{cajaalerta}[1][]{%
  enhanced, breakable, sharp corners,
  colback=acento!6, colframe=acento,
  boxrule=0.4pt, leftrule=3pt,
  fonttitle=\bfseries\color{white},
  coltitle=white, colbacktitle=acento,
  title={Importante}, #1
}

% ===================== Encabezado y pie ==========================
\pagestyle{fancy}
\fancyhf{}
\fancyhead[L]{\small\textcolor{primario}{\textbf{Programa de Capacitación en Desarrollo con IA}}}
\fancyhead[R]{\small\textcolor{primario}{\thepage}}
\fancyfoot[C]{\small\textcolor{textoGris}{Confidencial -- Uso interno}}
\renewcommand{\headrulewidth}{0.4pt}
\renewcommand{\footrulewidth}{0pt}

% ===================== Formato de títulos ========================
\usepackage{titlesec}
\titleformat{\section}{\Large\bfseries\color{primario}}{\thesection.}{0.6em}{}
\titleformat{\subsection}{\large\bfseries\color{secundario}}{\thesubsection}{0.5em}{}
\titleformat{\subsubsection}{\normalsize\bfseries\color{textoGris}}{}{0pt}{}

% ===================== Listas compactas ==========================
\usepackage{enumitem}
\setlist[itemize]{leftmargin=1.4em,itemsep=2pt,topsep=2pt}
\setlist[enumerate]{leftmargin=1.8em,itemsep=2pt,topsep=2pt}

% ===================== Encabezado de tabla =======================
\renewcommand{\arraystretch}{1.25}

% ===================== Metadatos del PDF =========================
\usepackage[hidelinks]{hyperref}
\usepackage{bookmark}
\hypersetup{
  pdftitle={Programa de Capacitación en Desarrollo con IA},
  pdfauthor={Equipo de Tecnología de la Información},
  pdfsubject={Programa formativo de 18 horas},
  pdfkeywords={IA, desarrollo, capacitación, TI},
  colorlinks=true, linkcolor=primario, urlcolor=secundario
}

% =================================================================
\begin{document}

% ============================ PORTADA ============================
\begin{titlepage}
  \thispagestyle{empty}
  \begin{center}
    \vspace*{1.5cm}
    \textcolor{primario}{\rule{\linewidth}{1.2pt}}\\[0.4cm]
    {\Huge\bfseries\textcolor{primario}{Programa de Capacitación}\\[0.3cm]
     Avanzada en Desarrollo con Inteligencia Artificial\\[0.3cm]
     para Equipos de TI\par}
    \vspace{0.4cm}
    \textcolor{primario}{\rule{\linewidth}{1.2pt}}\\[1.2cm]

    {\Large\color{textoGris}\itshape Ingeniería de Software, Análisis de Datos,\\
     Automatización y Gestión de Proyectos con IA\par}

    \vspace{2cm}
    \begin{tcolorbox}[enhanced, width=0.7\linewidth, colback=fondoClaro,
        colframe=secundario, boxrule=0.5pt, halign=center, sharp corners]
      \large
      \textbf{18 horas} de formación\\[0.2cm]
      \textbf{6 módulos} temáticos\\[0.2cm]
      \textbf{De los fundamentos a la producción}
    \end{tcolorbox}

    \vfill
    {\large\color{textoGris}%
     \textbf{Departamento de Tecnología de la Información}\\[0.2cm]
     \today\par}
    \vspace{1cm}
  \end{center}
\end{titlepage}

% ============================ ÍNDICE ==============================
\tableofcontents
\newpage

% =================================================================
\section{Presentación del programa}
% =================================================================

\begin{cajaobjetivo}
Este programa está diseñado para equipos de TI que quieren dominar la IA aplicada al desarrollo de software. A lo largo de \textbf{18 horas}, distribuidas en \textbf{6 módulos progresivos}, se transita desde los fundamentos matemáticos de la IA hasta su aplicación en producción, combinando teoría, laboratorios prácticos y debates técnicos.
\end{cajaobjetivo}

La incorporación de la inteligencia artificial en los flujos de trabajo de tecnología de la información ha dejado de ser un mero complemento de autocompletado para convertirse en un cambio fundamental en la arquitectura de sistemas. Este programa dota a los ingenieros de TI de las competencias necesarias para actuar como \textbf{directores y facilitadores} de esta transformación tecnológica.

\subsection{¿A quién está dirigido?}

\begin{itemize}
  \item Ingenieros de TI y desarrolladores de software.
  \item Analistas de datos y arquitectos de sistemas.
  \item Líderes técnicos y gestores de proyectos.
\end{itemize}

\begin{cajanota}
No es necesario ser experto en IA: el programa comienza desde cero y sube el nivel de forma gradual, alternando explicaciones sencillas con contenido técnico avanzado.
\end{cajanota}

% =================================================================
\section{Metodología del curso}
% =================================================================

El curso combina tres componentes pedagógicos en cada clase:

\begin{table}[h]
  \centering
  \renewcommand{\arraystretch}{1.3}
  \rowcolors{2}{fondoClaro}{white}
  \begin{tabular}{>{\bfseries\color{primario}}p{4.5cm} p{9.5cm}}
    \toprule
    \rowcolor{primario}\color{white}\textbf{Componente} & \color{white}\textbf{Qué incluye} \\
    \midrule
    Clase teórica & Conceptos explicados con ejemplos sencillos, analogías cotidianas y, progresivamente, profundidad técnica. \\
    Laboratorio práctico & Ejercicios guiados con IA ejecutándose en local, con apoyo del instructor en tiempo real. \\
    Preguntas de debate & Reflexión y discusión en grupo sobre casos reales. \\
    \bottomrule
  \end{tabular}
\end{table}

\begin{cajaalerta}
El docente evitará exposiciones teóricas superiores a \textbf{45 minutos} por sesión. La IA local debe estar activada en todo momento durante los laboratorios.
\end{cajaalerta}

% =================================================================
\section{Estructura general del programa}
% =================================================================

El programa se organiza en \textbf{seis módulos} acumulativos, cada uno de tres horas, garantizando un equilibrio óptimo entre fundamentación científica, demostración de patrones arquitectónicos y experimentación directa.

\begin{table}[h]
  \centering
  \small
  \renewcommand{\arraystretch}{1.35}
  \rowcolors{2}{fondoClaro}{white}
  \begin{tabularx}{\linewidth}{c Y Y Y}
    \toprule
    \rowcolor{primario}\color{white}\textbf{Módulo} & \color{white}\textbf{Horas} & \color{white}\textbf{Temas clave} & \color{white}\textbf{Herramientas} \\
    \midrule
    I. Fundamentos de IA y NLP & 1--3 & Teoría de sistemas, atención, MoE, LoRA y RoPE. & Visualizadores de atención y tokenización. \\
    II. Prompts y RAG & 4--6 & STAR, GOLD, pipelines RAG, GraphRAG, Dify. & Dify, Flowise, bases vectoriales. \\
    III. Agentes y SDD & 7--9 & MCP, ADE híbridos, Spec-Driven Development. & Claude Code, OpenCode, Gentle-AI. \\
    IV. Documentación y DDD & 10--12 & Markdown, Tolaria, Engram, DDD, contratos API. & Tolaria, Engram, generadores de contratos. \\
    V. Datos y Automatización & 13--15 & EDA automático, Data Formulator, n8n. & Sweetviz, YData-Profiler, n8n. \\
    VI. Pruebas, Gestión y Frontend & 16--18 & TDD/BDD, Playwright, AEPD/LCSP, OpenPencil. & Playwright, OpenPencil con MCP. \\
    \bottomrule
  \end{tabularx}
\end{table}

% =================================================================
\section{Detalle por módulo}
% =================================================================

\subsection{Módulo I: Fundamentos de IA y NLP Moderno (horas 1--3)}

Se aborda la IA desde la teoría de sistemas, desmitificando su naturaleza como ``entidad pensante'' para definirla como un sistema de aproximación de funciones. Se cubren tokenización, embeddings, la arquitectura Transformer, modelos autoregresivos, Mezcla de Expertos, LoRA y RoPE.

\paragraph{Clases:}
\begin{itemize}
  \item \textbf{Hora 1.} Teoría de sistemas, álgebra y aproximación universal.
  \item \textbf{Hora 2.} LLMs y procesamiento del lenguaje natural moderno.
  \item \textbf{Hora 3.} Arquitectura Transformer, sequence-to-sequence y escalamiento.
\end{itemize}

\subsection{Módulo II: Prompting Avanzado y Arquitecturas RAG (horas 4--6)}

Se trasciende el uso intuitivo de los prompts para estudiarlos como paradigma de programación declarativa. Se introduce el patrón Retrieval-Augmented Generation (RAG) y las plataformas de orquestación visual.

\paragraph{Clases:}
\begin{itemize}
  \item \textbf{Hora 4.} Estructuración metódica de prompts (STAR, RTF, GOLD, CoT/ToT).
  \item \textbf{Hora 5.} RAG, embeddings vectoriales y GraphRAG.
  \item \textbf{Hora 6.} Orquestadores: LangChain, LangGraph, Dify y Flowise.
\end{itemize}

\subsection{Módulo III: Agentes Inteligentes y Entornos ADE (horas 7--9)}

Se detalla cómo los agentes interactúan con el entorno mediante el protocolo MCP, se comparan IDEs híbridos frente a CLIs autónomos, y se presenta la metodología Spec-Driven Development (SDD).

\paragraph{Clases:}
\begin{itemize}
  \item \textbf{Hora 7.} Anatomía de un agente y el protocolo MCP.
  \item \textbf{Hora 8.} Entornos de desarrollo agénticos: Cursor, Claude Code, OpenCode.
  \item \textbf{Hora 9.} Spec-Driven Development con la suite Gentleman-AI.
\end{itemize}

\subsection{Módulo IV: Documentación, Memoria y DDD (horas 10--12)}

Se explora la documentación Markdown como interfaz bidireccional humano-IA, los sistemas de memoria persistente para mitigar la amnesia contextual, y la arquitectura Domain-Driven Design con contratos de APIs.

\paragraph{Clases:}
\begin{itemize}
  \item \textbf{Hora 10.} Documentación técnica y bases de conocimiento (Tolaria, Obsidian).
  \item \textbf{Hora 11.} Memoria persistente: Engram y Codebase-Memory-MCP.
  \item \textbf{Hora 12.} Arquitectura DDD, contratos Contract-First y JSON-First.
\end{itemize}

\subsection{Módulo V: Análisis de Datos y Automatización con n8n (horas 13--15)}

Se capacita al equipo en la optimización de flujos analíticos descriptivos mediante EDA automatizado, la visualización guiada por IA con Data Formulator, y la orquestación de flujos reactivos con n8n.

\paragraph{Clases:}
\begin{itemize}
  \item \textbf{Hora 13.} Análisis exploratorio de datos automatizado (Sweetviz, YData-Profiler).
  \item \textbf{Hora 14.} Data Formulator de Microsoft Research.
  \item \textbf{Hora 15.} Automatización de flujos de trabajo con n8n.
\end{itemize}

\subsection{Módulo VI: Pruebas, Gestión y Frontend (horas 16--18)}

Se aborda el aseguramiento de calidad moderno (TDD/BDD, Playwright), la gestión de proyectos asistida por IA con cumplimiento normativo (AEPD/LCSP), y la generación de frontend con OpenPencil conectado vía MCP.

\paragraph{Clases:}
\begin{itemize}
  \item \textbf{Hora 16.} Aseguramiento de calidad, pirámide de pruebas e IA.
  \item \textbf{Hora 17.} Gestión de proyectos de TI asistida por IA.
  \item \textbf{Hora 18.} Generación de interfaces frontend con OpenPencil y MCP.
\end{itemize}

% =================================================================
\section{Herramientas que se dominarán}
% =================================================================

\begin{table}[h]
  \centering
  \small
  \renewcommand{\arraystretch}{1.3}
  \rowcolors{2}{fondoClaro}{white}
  \begin{tabularx}{\linewidth}{Y Y}
    \toprule
    \rowcolor{primario}\color{white}\textbf{Categoría} & \color{white}\textbf{Herramientas} \\
    \midrule
    Modelos de lenguaje & GPT, LLaMA, BERT, TinyLlama, Phi-2. \\
    Agentes y ADE & Claude Code, OpenCode, Codex, Gemini CLI. \\
    Metodología SDD & Gentle-AI, gentle-pi, Engram. \\
    Protocolos & MCP (Model Context Protocol). \\
    RAG y orquestación & LangChain, LangGraph, Dify, Flowise, ChromaDB. \\
    Análisis de datos & Sweetviz, YData-Profiler, Data Formulator. \\
    Automatización & n8n. \\
    Pruebas & Playwright Codegen, Gherkin. \\
    Documentación & Tolaria, Obsidian, Markdown + YAML frontmatter. \\
    Frontend & OpenPencil con servidor MCP. \\
    \bottomrule
  \end{tabularx}
\end{table}

\begin{cajanota}
Todas las herramientas del programa son de código abierto o disponen de una versión gratuita ejecutable en local, garantizando privacidad e independencia tecnológica.
\end{cajanota}

% =================================================================
\section{Resultados de aprendizaje esperados}
% =================================================================

Al finalizar el programa, los participantes serán capaces de:

\begin{enumerate}
  \item \textbf{Comprender} los fundamentos matemáticos y técnicos de la IA moderna, sin mitos ni confusiones.
  \item \textbf{Comunicarse} con la IA mediante prompts estructurados y precisos (STAR, GOLD, CoT).
  \item \textbf{Implementar} arquitecturas RAG y GraphRAG sobre conocimiento corporativo cerrado.
  \item \textbf{Usar} agentes autónomos en entornos de desarrollo reales con el protocolo MCP.
  \item \textbf{Estructurar} el conocimiento del equipo mediante documentación Markdown y memoria persistente.
  \item \textbf{Diseñar} arquitecturas Domain-Driven Design con contratos de APIs Contract-First.
  \item \textbf{Automatizar} análisis de datos exploratorios y flujos de trabajo reactivos.
  \item \textbf{Asegurar} la calidad del código generado por IA mediante TDD/BDD y Playwright.
  \item \textbf{Aplicar} el marco legal AEPD/LCSP en la gestión de proyectos con IA.
\end{enumerate}

% =================================================================
\section{Lineamientos para la ejecución}
% =================================================================

\subsection{Enfoque de aprendizaje práctico continuo}

Cada hora lectiva especifica su práctica de laboratorio. Es mandatorio que el docente guíe, en tiempo real, la ejecución de comandos y la construcción de prompts correctos utilizando el proyector del aula.

\subsection{Uso de la tercera persona en el material}

Para mantener la formalidad corporativa y el rigor académico, todo el material visual se redacta en voz pasiva impersonal del español (``se define'', ``se ejecuta'', ``se realiza'').

\subsection{Gobierno técnico y gestión de errores}

Durante los laboratorios es natural que las llamadas a las APIs devuelvan excepciones de red, límites de velocidad o código erróneo. El docente debe aprovechar estas fallas de forma pedagógica para enseñar Spec-Driven Development, forzando a los agentes a leer los logs, consultar la memoria local de Engram y aplicar parches de corrección autónomos.

\begin{cajaalerta}
La IA es un colaborador analítico y de ejecución. El profesional de TI actúa como \textbf{director y árbitro} de las discusiones, asegurando el cumplimiento de las restricciones técnicas en cada etapa del ciclo de desarrollo.
\end{cajaalerta}

% =================================================================
\section{Cronograma resumido}
% =================================================================

\begin{table}[h]
  \centering
  \renewcommand{\arraystretch}{1.3}
  \rowcolors{2}{fondoClaro}{white}
  \begin{tabular}{>{\bfseries\color{primario}}c p{8cm} c}
    \toprule
    \rowcolor{primario}\color{white}\textbf{Bloque} & \color{white}\textbf{Módulo} & \color{white}\textbf{Duración} \\
    \midrule
    Bloque 1 & Fundamentos de IA y NLP Moderno & 3 horas \\
    Bloque 2 & Prompting Avanzado y Arquitecturas RAG & 3 horas \\
    Bloque 3 & Agentes Inteligentes y Entornos ADE & 3 horas \\
    Bloque 4 & Documentación, Memoria y DDD & 3 horas \\
    Bloque 5 & Análisis de Datos y Automatización con n8n & 3 horas \\
    Bloque 6 & Pruebas, Gestión y Frontend & 3 horas \\
    \midrule
    \textbf{Total} & \textbf{Programa completo} & \textbf{18 horas} \\
    \bottomrule
  \end{tabular}
\end{table}

\vspace{0.5cm}
\begin{center}
  \begin{tcolorbox}[enhanced, width=0.85\linewidth, colback=verde!8, colframe=verde,
      boxrule=0.5pt, halign=center, sharp corners]
    \large\textbf{El equipo de TI está listo para liderar la transformación con IA.}
  \end{tcolorbox}
\end{center}

\vfill
\begin{center}
  \small\color{textoGris}%
  \emph{Documento generado como programa oficial del curso.}\\
  Departamento de Tecnología de la Información -- \today
\end{center}

\end{document}
